\documentclass[12pt,a4paper]{exam}
\usepackage[utf8]{inputenc}
\usepackage{amsmath,amssymb,amsfonts}
\usepackage{geometry}
\usepackage{enumitem}
\usepackage{graphicx}
\usepackage{hyperref}
\usepackage{xcolor}
\geometry{a4paper, top=2cm, bottom=2cm, left=2.5cm, right=2.5cm}
\hypersetup{colorlinks=true, linkcolor=blue, urlcolor=blue}
\renewcommand{\thequestion}{\textbf{\arabic{question}}}
\renewcommand{\questionlabel}{\thequestion.}
\pointname{ marks}
\pointsdroppedatright
\bracketedpoints
\setlength{\parindent}{0pt}
\setlength{\emergencystretch}{3em}
\allowdisplaybreaks
\newsavebox{\schmathbox}
\newcommand{\fitmath}[1]{%
  \sbox{\schmathbox}{$\displaystyle #1$}%
  \ifdim\wd\schmathbox>\linewidth
    \resizebox{\linewidth}{!}{\usebox{\schmathbox}}%
  \else
    \usebox{\schmathbox}%
  \fi}

\begin{document}

\thispagestyle{empty}
\begin{center}
  \textbf{\LARGE Diag} \\[0.3cm]
  \textbf{Time Allowed: 3 Hours} \hfill \textbf{Maximum Marks: 80} \\[0.3cm]
  \rule{\textwidth}{0.5pt}
\end{center}

\vspace{0.3cm}

\noindent\textbf{General Instructions:}
\begin{enumerate}
  \item {\normalfont\normalcolor This question paper contains 40 questions. All questions are compulsory.}
  \item {\normalfont\normalcolor Questions are grouped into sections A through E.}
  \item {\normalfont\normalcolor Use of calculators is NOT permitted.}
\end{enumerate}
\bigskip
\noindent\textbf{\large Section A: Multiple Choice \& Assertion-Reason (1 mark each)}

\begin{questions}
\question[1] The value of $(1 + \tan^2 \theta) \cos^2 \theta$ is
\begin{enumerate}[label=(\Alph*)]
  \item (A) 1
  \item (B) 0
  \item (C) 2
  \item (D) -1
\end{enumerate}
\vspace{0.15cm}

\question[1] If $\sin A = \cos A$, then the value of $A$ is
\begin{enumerate}[label=(\Alph*)]
  \item $(A) 30^\circ$
  \item $(B) 45^\circ$
  \item $(C) 60^\circ$
  \item $(D) 0^\circ$
\end{enumerate}
\vspace{0.15cm}

\question[1] The value of $\text{cosec}^2 \theta - 1$ is equal to
\begin{enumerate}[label=(\Alph*)]
  \item $(A) \tan^2 \theta$
  \item $(B) \sin^2 \theta$
  \item $(C) \cos^2 \theta$
  \item $(D) \cot^2 \theta$
\end{enumerate}
\vspace{0.15cm}

\question[1] The 4th term of an AP is 14 and the 12th term is 70. What is the first term $a$?
\begin{enumerate}[label=(\Alph*)]
  \item (A) -7
  \item (B) 7
  \item (C) -14
  \item (D) 14
\end{enumerate}
\vspace{0.15cm}

\question[1] If the sum of first $n$ terms of an AP is $S_n = 3n^2 + n$, what is the common difference $d$?
\begin{enumerate}[label=(\Alph*)]
  \item (A) 3
  \item (B) 4
  \item (C) 6
  \item (D) 2
\end{enumerate}
\vspace{0.15cm}

\question[1] Which term of the AP: 21, 18, 15, ... is -81?
\begin{enumerate}[label=(\Alph*)]
  \item (A) 34th
  \item (B) 35th
  \item (C) 36th
  \item (D) 37th
\end{enumerate}
\vspace{0.15cm}

\question[1] In an AP, if $a=5$ and $d=3$, find the sum of the first 10 terms ($S_{10}$)?
\begin{enumerate}[label=(\Alph*)]
  \item (A) 150
  \item (B) 175
  \item (C) 185
  \item (D) 195
\end{enumerate}
\vspace{0.15cm}

\question[1] If $k, 2k-1, 2k+1$ are three consecutive terms of an AP, find the value of $k$.
\begin{enumerate}[label=(\Alph*)]
  \item (A) 3
  \item (B) 2
  \item (C) 4
  \item (D) 5
\end{enumerate}
\vspace{0.15cm}

\question[1] Find the sum of all odd numbers between 0 and 50.
\begin{enumerate}[label=(\Alph*)]
  \item (A) 600
  \item (B) 625
  \item (C) 650
  \item (D) 675
\end{enumerate}
\vspace{0.15cm}

\question[1] The 10th term from the end of the AP: 5, 8, 11, ..., 185 is:
\begin{enumerate}[label=(\Alph*)]
  \item (A) 150
  \item (B) 152
  \item (C) 155
  \item (D) 158
\end{enumerate}
\vspace{0.15cm}

\question[1] If the sum of first $n$ terms is $n^2$, then the $n$th term is:
\begin{enumerate}[label=(\Alph*)]
  \item (A) 2n+1
  \item (B) 2n
  \item (C) 2n-1
  \item (D) n-1
\end{enumerate}
\vspace{0.15cm}

\question[1] How many two-digit numbers are divisible by 3?
\begin{enumerate}[label=(\Alph*)]
  \item (A) 30
  \item (B) 31
  \item (C) 32
  \item (D) 33
\end{enumerate}
\vspace{0.15cm}

\question[1] What is the common difference of an AP in which $a_{25} - a_{20} = 45$?
\begin{enumerate}[label=(\Alph*)]
  \item (A) 5
  \item (B) 9
  \item (C) 10
  \item (D) 4.5
\end{enumerate}
\vspace{0.15cm}

\question[1] Assertion (A): For an angle $\theta$, $\sin^2 \theta + \cos^2 \theta = 1$. Reason (R): This identity holds true for all values of $\theta$ where $0^\circ \le \theta \le 90^\circ$.
\begin{enumerate}[label=(\Alph*)]
  \item (A) Both Assertion (A) and Reason (R) are true and Reason (R) is the correct explanation of Assertion (A)
  \item (B) Both Assertion (A) and Reason (R) are true but Reason (R) is NOT the correct explanation of Assertion (A)
  \item (C) Assertion (A) is true but Reason (R) is false
  \item (D) Assertion (A) is false but Reason (R) is true
\end{enumerate}
\vspace{0.15cm}

\question[1] Assertion (A): The value of $\tan 45^\circ$ is $1$. Reason (R): $\tan \theta = \frac{\sin \theta}{\cos \theta}$ and $\sin 45^\circ = \cos 45^\circ = \frac{1}{\sqrt{2}}$.
\begin{enumerate}[label=(\Alph*)]
  \item (A) Both Assertion (A) and Reason (R) are true and Reason (R) is the correct explanation of Assertion (A)
  \item (B) Both Assertion (A) and Reason (R) are true but Reason (R) is NOT the correct explanation of Assertion (A)
  \item (C) Assertion (A) is true but Reason (R) is false
  \item (D) Assertion (A) is false but Reason (R) is true
\end{enumerate}
\vspace{0.15cm}

\question[1] Assertion (A): The value of $\sin \theta$ can be $1.5$ for some angle $\theta$. Reason (R): The range of $\sin \theta$ is $[-1, 1]$.
\begin{enumerate}[label=(\Alph*)]
  \item (A) Both Assertion (A) and Reason (R) are true and Reason (R) is the correct explanation of Assertion (A)
  \item (B) Both Assertion (A) and Reason (R) are true but Reason (R) is NOT the correct explanation of Assertion (A)
  \item (C) Assertion (A) is true but Reason (R) is false
  \item (D) Assertion (A) is false but Reason (R) is true
\end{enumerate}
\vspace{0.15cm}

\question[1] Assertion (A): $\sec^2 \theta - \tan^2 \theta = 1$. Reason (R): $1 + \tan^2 \theta = \sec^2 \theta$.
\begin{enumerate}[label=(\Alph*)]
  \item (A) Both Assertion (A) and Reason (R) are true and Reason (R) is the correct explanation of Assertion (A)
  \item (B) Both Assertion (A) and Reason (R) are true but Reason (R) is NOT the correct explanation of Assertion (A)
  \item (C) Assertion (A) is true but Reason (R) is false
  \item (D) Assertion (A) is false but Reason (R) is true
\end{enumerate}
\vspace{0.15cm}

\question[1] Assertion (A): If $\cos A = \frac{1}{2}$, then $\sin A = \frac{\sqrt{3}}{2}$. Reason (R): $\sin^2 A + \cos^2 A = 1$.
\begin{enumerate}[label=(\Alph*)]
  \item (A) Both Assertion (A) and Reason (R) are true and Reason (R) is the correct explanation of Assertion (A)
  \item (B) Both Assertion (A) and Reason (R) are true but Reason (R) is NOT the correct explanation of Assertion (A)
  \item (C) Assertion (A) is true but Reason (R) is false
  \item (D) Assertion (A) is false but Reason (R) is true
\end{enumerate}
\vspace{0.15cm}

\end{questions}
\bigskip
\noindent\textbf{\large Section B: Very Short Answer (2 marks each)}

\begin{questions}
\question[2] If $\sin A = \cos A$, where $0 < A < 90^\circ$, find the value of $A$.
\vspace{0.15cm}

\question[2] Evaluate $\frac{\tan 65^\circ}{\cot 25^\circ}$.
\vspace{0.15cm}

\question[2] Given $15 \cot A = 8$, find the value of $\sin A$.
\vspace{0.15cm}

\question[2] Find the value of $2 \tan^2 45^\circ + \cos^2 30^\circ - \sin^2 60^\circ$.
\vspace{0.15cm}

\question[2] Prove that $\sec^2 \theta - 1 = \tan^2 \theta$.
\vspace{0.15cm}

\question[2] Find the $11^{th}$ term of an Arithmetic Progression (AP) whose first term is $-3$ and the common difference is $4$.
\vspace{0.15cm}

\question[2] Which term of the AP $21, 18, 15, \dots$ is $-81$?
\vspace{0.15cm}

\question[2] Find the sum of the first $10$ multiples of $5$.
\vspace{0.15cm}

\question[2] If the $n^{th}$ term of an AP is given by $a_n = 3 + 2n$, find its common difference.
\vspace{0.15cm}

\question[2] Find the value of $k$ for which the terms $k-1, k+3, 3k-1$ are in Arithmetic Progression.
\vspace{0.15cm}

\end{questions}
\bigskip
\noindent\textbf{\large Section C: Short Answer (3 marks each)}

\begin{questions}
\question[3] If $\sqrt{3} \tan \theta = 1$, find the value of $\sin^2 \theta - \cos^2 \theta$.
\vspace{0.15cm}

\question[3] Prove that $\frac{\cos A}{1 + \sin A} + \frac{1 + \sin A}{\cos A} = 2 \sec A$.
\vspace{0.15cm}

\question[3] If $\tan A = \frac{3}{4}$, find all other trigonometric ratios for angle $A$.
\vspace{0.15cm}

\question[3] Find the sum of all two-digit numbers which, when divided by 4, leave 1 as a remainder.
\vspace{0.15cm}

\question[3] If the sum of the first $n$ terms of an AP is $S_n = 3n^2 + 5n$, find the 16th term of the AP.
\vspace{0.15cm}

\question[3] The 4th term of an AP is 0. Prove that the 25th term of the AP is three times its 11th term.
\vspace{0.15cm}

\question[3] How many terms of the AP 24, 21, 18, ... must be taken so that their sum is 78?
\vspace{0.15cm}

\question[3] Find the middle term of the Arithmetic Progression 7, 13, 19, ..., 241.
\vspace{0.15cm}

\end{questions}
\bigskip
\noindent\textbf{\large Section D: Long Answer (4-5 marks each)}

\begin{questions}
\question[5] Prove that $\frac{\tan \theta}{1 - \cot \theta} + \frac{\cot \theta}{1 - \tan \theta} = 1 + \sec \theta \csc \theta$.
\vspace{0.15cm}

\question[5] If $\sec \theta + \tan \theta = p$, prove that $\sin \theta = \frac{p^2 - 1}{p^2 + 1}$.
\vspace{0.15cm}

\question[5] The sum of the first $n$ terms of an Arithmetic Progression is given by $S_n = 3n^2 + 5n$. Find the $n^{th}$ term of this AP and hence find the $20^{th}$ term.
\vspace{0.15cm}

\question[5] In an AP, the sum of first 10 terms is $-150$ and the sum of its next 10 terms is $-550$. Find the Arithmetic Progression.
\vspace{0.15cm}

\end{questions}
\bigskip

\end{document}
